% Generated draft: requires mathematical and source review.
\documentclass[11pt]{article}
\usepackage[margin=25mm]{geometry}
\usepackage{fontspec}
\setmainfont{Latin Modern Roman}
\usepackage{amsmath,amssymb,mathtools,unicode-math}
\setmathfont{Latin Modern Math}
\usepackage{enumitem,needspace,xurl,hyperref}
\setlength{\parindent}{0pt}
\setlength{\parskip}{5pt}
\setlength{\emergencystretch}{3em}
\newcommand{\N}{\mathbb{N}}
\newcommand{\Z}{\mathbb{Z}}
\newcommand{\Q}{\mathbb{Q}}
\newcommand{\R}{\mathbb{R}}
\newcommand{\C}{\mathbb{C}}
\newcommand{\F}{\mathbb{F}}
\newcommand{\A}{\mathbb{A}}
\newcommand{\PP}{\mathbb{P}}
\newcommand{\E}{\mathbb{E}}
\newcommand{\eps}{\varepsilon}
\newcommand{\dd}{\,\mathrm{d}}
\begin{document}
\section*{UnsolvedMath ID: 6400009}
\textbf{Source problem code:} catalogue entry
\subsection{Definitions and mathematical statement}
A calibration on a Riemannian manifold is a closed \(k\)-form \(\alpha\) with comass at most \(1\). An oriented \(k\)-submanifold \(L\) is calibrated when \(\alpha|_L\) is its volume form. The source asks for compactness and singularity theories sufficient for enumerative invariants in three standard calibrated settings:
\begin{itemize}
\item In a Calabi--Yau manifold of complex dimension \(m\), with K\"ahler form \(\omega\) and normalized holomorphic volume form \(\Omega\), a special Lagrangian \(m\)-submanifold of phase \(\theta\) satisfies \(\omega|_L=0\), \(\operatorname{Im}(e^{-i\theta}\Omega)|_L=0\), and is oriented so that \(\operatorname{Re}(e^{-i\theta}\Omega)|_L\) is its volume form.
\item In a torsion-free \(G_2\) seven-manifold with calibrating three-form \(\varphi\), an associative three-submanifold is calibrated by \(\varphi\).
\item In the same setting, a co-associative four-submanifold is calibrated by \(*\varphi\).
\end{itemize}
Develop compactness and singularity descriptions of the appropriate moduli spaces, including limiting multiplicities and multiple-cover phenomena, that permit well-defined enumerative invariants. This retains the source as a research program rather than claiming a single compactness theorem would itself define a count.

The excerpt does not prescribe compact versus noncompact ambient spaces, boundary/asymptotic conditions, fixed homology classes, varying geometric structures, the admissible singular objects and convergence, or the orientation, weighting, and deformation-invariance requirements of the proposed invariant. Those data are necessary to state a particular counting theory and are not guessed here.

The supplied summary reports compactness-with-multiplicity and what it calls ``countable-compactness'' results for calibrated submanifolds, while saying singularity and multiple-cover theory does not yet supply robust invariants in full generality. It provides neither exact theorem statements nor a definition of that compactness terminology. This is a qualified partial-progress assessment, distinct from independent verification of the 2008 program's present status.
\subsection{Short English statement}
Build compactness, singularity, and multiple-cover theories for special Lagrangian, associative, and co-associative submanifolds that support meaningful enumerative counts. The source does not specify a particular moduli problem or invariant, so the program remains broader than one precise theorem.
\subsection{Sources}
\begin{itemize}
\item Local registry: \url{data/unsolvedmath/problems.json}, permanent UnsolvedMath ID \(6400009\), source code catalogue entry; supplied statement, background, and research summary, updated 2026-08-06.
\item Simon Donaldson, \emph{Some Problems in Differential Geometry and Topology} (2008), explicit question/program 9: \url{https://wwwf.imperial.ac.uk/~skdona/NONLINEARITYPROBLEMS.PDF}. Supplied normalized PDF-text extraction accessed 2026-07-29; current status marked as needing review.
\item The summary reports multiplicity and compactness progress without separate original-paper references or exact definitions.
\end{itemize}
\end{document}
